% Use case: noise-aware, post-selected state preparation in the XY register (continues Section sec:noisy).
% Numbers: scripts/noisy_xy_control.py -> generated/noisy_control.tex (macros \NC...), figures/noisy_control.pdf.
% The file loads its macros itself if the preamble has not (\gen{noisy_control}).
\IfFileExists{generated/noisy_control.tex}{\ifdefined\NCRows\else% generated by scripts/noisy_xy_control.py -- do not edit by hand
\newcommand{\NCRows}{%
4 & (i) & 2.73 & 1.0000 & 0.941 & 0.929 & 0.987 & 0.795 & 0.948\\
4 & (ii) & 2.66 & 0.9993 & 0.942 & 0.930 & 0.988 & 0.801 & 0.951\\
4 & (iii) & 2.71 & 1.0000 & 0.941 & 0.930 & 0.988 & 0.799 & 0.952\\
6 & (i) & 5.49 & 0.9981 & 0.846 & 0.807 & 0.952 & 0.490 & 0.816\\
6 & (ii) & 4.88 & 0.9889 & 0.860 & 0.821 & 0.953 & 0.535 & 0.846\\
6 & (iii) & 5.14 & 0.9949 & 0.854 & 0.818 & 0.956 & 0.519 & 0.840\\
}
\newcommand{\NCFourMinFPostOne}{0.948}
\newcommand{\NCFourMinFPostTwo}{0.951}
\newcommand{\NCFourMinFPostThree}{0.952}
\newcommand{\NCFourTOne}{2.73}
\newcommand{\NCFourTTwo}{2.66}
\newcommand{\NCFourTThree}{2.71}
\newcommand{\NCFourFPostOne}{0.987}
\newcommand{\NCFourFPostTwo}{0.988}
\newcommand{\NCFourFPostThree}{0.988}
\newcommand{\NCFourFOne}{0.929}
\newcommand{\NCFourFTwo}{0.930}
\newcommand{\NCFourFThree}{0.930}
\newcommand{\NCFourPOne}{0.941}
\newcommand{\NCFourPTwo}{0.942}
\newcommand{\NCFourPThree}{0.941}
\newcommand{\NCFourFZeroOne}{1.0000}
\newcommand{\NCFourFZeroTwo}{0.9993}
\newcommand{\NCFourFZeroThree}{1.0000}
\newcommand{\NCFourGainPost}{1.07}
\newcommand{\NCFourGainRaw}{1.02}
\newcommand{\NCFourGainTotal}{5.89}
\newcommand{\NCFourPostFactor}{5.6}
\newcommand{\NCFourRatioLow}{0.94}
\newcommand{\NCFourRatioHigh}{0.94}
\newcommand{\NCFourSolutions}{32}
\newcommand{\NCFourSeconds}{130}
\newcommand{\NCSixMinFPostOne}{0.816}
\newcommand{\NCSixMinFPostTwo}{0.846}
\newcommand{\NCSixMinFPostThree}{0.840}
\newcommand{\NCSixTOne}{5.49}
\newcommand{\NCSixTTwo}{4.88}
\newcommand{\NCSixTThree}{5.14}
\newcommand{\NCSixFPostOne}{0.952}
\newcommand{\NCSixFPostTwo}{0.953}
\newcommand{\NCSixFPostThree}{0.956}
\newcommand{\NCSixFOne}{0.807}
\newcommand{\NCSixFTwo}{0.821}
\newcommand{\NCSixFThree}{0.818}
\newcommand{\NCSixPOne}{0.846}
\newcommand{\NCSixPTwo}{0.860}
\newcommand{\NCSixPThree}{0.854}
\newcommand{\NCSixFZeroOne}{0.9981}
\newcommand{\NCSixFZeroTwo}{0.9889}
\newcommand{\NCSixFZeroThree}{0.9949}
\newcommand{\NCSixGainPost}{1.09}
\newcommand{\NCSixGainRaw}{1.07}
\newcommand{\NCSixGainTotal}{4.07}
\newcommand{\NCSixPostFactor}{4.0}
\newcommand{\NCSixRatioLow}{1.29}
\newcommand{\NCSixRatioHigh}{0.90}
\newcommand{\NCSixSolutions}{32}
\newcommand{\NCSixSeconds}{528}
\newcommand{\NCSegments}{4}
\newcommand{\NCSubsteps}{2}
\newcommand{\NCEvalSubsteps}{20}
\newcommand{\NCStarts}{32}
\newcommand{\NCWarm}{8}
\newcommand{\NCWarmLarge}{4}
\newcommand{\NCSettings}{12}
\newcommand{\NCSteps}{150}
\newcommand{\NCStepsZero}{400}
\newcommand{\NCFZeroMin}{0.99}
\newcommand{\NCLambda}{0.05}
\newcommand{\NCRecheck}{6\times10^{-4}}
\newcommand{\NCBirthDeath}{4\times10^{-6}}
\newcommand{\NCExactDiff}{6\times10^{-6}}
\newcommand{\NCGradErr}{6\times10^{-8}}
\newcommand{\NCSeconds}{658}
\newcommand{\NCMinutes}{11}
\newcommand{\NCThreads}{4}
\newcommand{\NCnMax}{6}
\fi}{}
\providecommand{\NCSegments}{4}\providecommand{\NCSettings}{12}\providecommand{\NCSteps}{--}\providecommand{\NCStepsZero}{--}\providecommand{\NCSubsteps}{2}
\providecommand{\NCStarts}{--}\providecommand{\NCWarm}{--}\providecommand{\NCWarmLarge}{--}\providecommand{\NCLambda}{0.05}\providecommand{\NCEvalSubsteps}{20}\providecommand{\NCnMax}{6}
\providecommand{\NCRows}{\multicolumn{9}{c}{\emph{run scripts/noisy\_xy\_control.py}}\\}
\subsection{Noise-aware state preparation with post-selection}\label{sec:control}

Section~\ref{sec:noisy} evaluated a fixed evolution. In practice the evolution is designed, and the design can take
into account the noise and the fact that the Hamming weight will be measured. XY-QAOA needs an initial state in the
feasible subspace, and the Dicke state $|D^n_w\rangle$, the uniform superposition of all strings of weight $w$, is the
standard choice~\cite{hadfield2019,wang2020}. Streif et al.~\cite{streif2021} computed how likely such circuits are
to stay in the feasible subspace under local noise. Botelho et al.~\cite{botelho2022} proposed post-selection on
Hamming-weight subspaces for QAOA, an instance of symmetry verification~\cite{cai2023}. This raises a design question. Does a preparation that will be
post-selected need different controls from one that will not, and how much does designing for the noise gain over
a noiseless design? The question is an optimisation over the controls of a Lindblad evolution in
$\M=\bigoplus_NM_{\binom nN}$, and its objective uses the centre of $\M$. It needs gradients through open-system
dynamics in the sectors, the central projection $P_w$, and many noise settings at once. These are the three things the
library provides.

\paragraph{Parametrisation.} A ring of $n$ qubits starts in $|1^w0^{n-w}\rangle$, $w=n/2$, and is driven during a time
$T$ by $S=\NCSegments$ piecewise-constant segments of length $T/S$ with Hamiltonians
\begin{equation}\label{eq:control}
H_s=\sum_{\langle ij\rangle}J_{ij,s}\,\tfrac12\bigl(S^+_iS^-_j+S^-_iS^+_j\bigr)+\sum_ih_{i,s}S^z_i
+V_s\sum_{\langle ij\rangle}S^z_iS^z_j,\qquad |J_{ij,s}|,|h_{i,s}|,|V_s|\le1,\quad 0<T\le2n.
\end{equation}
These are bond-resolved XY couplings, local fields and a uniform phase-separator term. Every term commutes with $N$,
so the controls stay in $\M$. The bond terms are \texttt{SpinChain.bond(i, j, Jxy, Jz)}, and $H_s$ is a linear
combination of them with trainable coefficients. Without the fields the target cannot be reached. The global spin
flip $\prod_iX_i$ commutes with all bond terms and maps sector $n/2$ to itself, the Dicke state is even under it, and
$|1^w0^{n-w}\rangle$ has equal weight on both parities, so $F\le\frac12$. With $V_s=0$ as well, the evolution is
quadratic in Jordan--Wigner fermions and cannot leave the Slater determinants. The noise consists of the four models of
Section~\ref{sec:noisy} ($T_1$ at zero temperature, thermal, dephasing and mixed), each at the rates $\gamma=0.003$,
$0.01$ and $0.03$, so the batch has $K=\NCSettings$ settings. Each segment is integrated over unit time with $H_s$
replaced by $(T/S)H_s$ and jump coefficients $\sqrt{\gamma T/S}$ (\texttt{SectorJump.scaled}), so $\rho(T)$ is
differentiable in all controls and in $T$.

\paragraph{Objectives.} With $F_k=\langle D^n_w|\rho_k(T)|D^n_w\rangle$ for noise setting $k$,
$p_k=\Tr(P_w\rho_k(T))$ the acceptance and $F_{\mathrm{post},k}=F_k/p_k$ the fidelity of the Lüders-conditioned state
(\texttt{condition\_on}), we compare three designs:
(i) the noiseless design, which maximises $F_0-\lambda T/T_{\max}$ without noise, $\lambda=\NCLambda$. The penalty
pushes $T$ close to the shortest time at which the controls reach $F_0\approx1$, which is the most favourable choice
for a design that ignores the noise;
(ii) the noise-aware design without post-selection, which maximises $\bar F=\frac1K\sum_kF_k$; and
(iii) the noise-aware design with post-selection, which maximises $\bar F_{\mathrm{post}}=\frac1K\sum_kF_{\mathrm{post},k}$.
By Proposition~\ref{prop:noisy}(a), $p_k$ depends on $T$ and the rates but not on the controls. At fixed $T$,
therefore, $F_k=p_k(T)F_{\mathrm{post},k}$, and (ii) and (iii) differ only in the weights $p_k$ of the settings and in
the choice of $T$: design (ii) pays for the detected loss $\Gamma_wT$, design (iii) does not. Each objective is
maximised by \texttt{optimize.extremize} (Adam, \NCSteps\ steps; \NCStepsZero\ for (i)). Every evaluation is a
single batched \texttt{lindblad\_evolve} per segment for all starts and all $K$ settings (\texttt{method='kraus'},
\NCSubsteps\ steps per segment). The noiseless design is the best of \NCStarts\ random starts. Designs (ii) and (iii)
start from the \NCWarm\ ($n=4$) or \NCWarmLarge\ ($n=6$) best noiseless solutions, and the starts of (iii) also include
the solution of (ii).

\begin{table}[t]
\centering\footnotesize
\setlength{\tabcolsep}{3.2pt}
\begin{tabular}{rlrrrrrrr}
\toprule
$n$ & design & $T$ & $F_0$ & $\bar p$ & $\bar F$ & $\bar F_{\mathrm{post}}$ & $\min_kF_k$ &
$\min_kF_{\mathrm{post},k}$\\
\midrule
\NCRows
\bottomrule
\end{tabular}
\caption[]{Dicke-state preparation on a ring of $n$ qubits by the three designs, evaluated over the $\NCSettings$ noise
settings (four models, $\gamma\in\{0.003,0.01,0.03\}$). $F_0$: fidelity without noise. Means and minima over the
settings; every value is from an independent evaluation with $\NCEvalSubsteps$ steps per segment. Full run, CPU
(4 threads).}
\label{tab:control}
\end{table}

\begin{figure}[t]
\centering\fig{noisy_control.pdf}{\textwidth}
\caption{Designs (i)--(iii) for $n=\NCnMax$ under mixed noise
($\gamma_\downarrow=\gamma_\varphi=\gamma$, $\gamma_\uparrow=0.1\gamma$) at rates inside and outside the training batch
(grey bars). (a) Infidelity with $|D^n_w\rangle$ with (solid) and without (dashed) post-selection. (b) The same,
relative to the noiseless design (i).}
\label{fig:control}
\end{figure}

\paragraph{Results.} Table~\ref{tab:control} and Figure~\ref{fig:control} show the outcome. For $n=6$ the three
designs use the times $T=\G{noisy_control}{\NCSixTOne}{--}$, $\G{noisy_control}{\NCSixTTwo}{--}$ and
$\G{noisy_control}{\NCSixTThree}{--}$. Design (ii) is the fastest and gives up the most noiseless fidelity
($F_0=\G{noisy_control}{\NCSixFZeroTwo}{--}$ against $\G{noisy_control}{\NCSixFZeroOne}{--}$ for (i)), because every
unit of time costs it the detected loss as well. Design (iii) ignores that loss, so it spends more time than (ii) but
still less than (i), because the undetected part of the noise (dephasing, and pairs of an up and a down jump) also
grows with $T$. The order of the times is
the same for $n=4$. Each noise-aware design wins on its own objective. Design (ii) has the largest mean raw fidelity,
$\bar F=\G{noisy_control}{\NCSixFTwo}{--}$ against $\G{noisy_control}{\NCSixFOne}{--}$ for (i), and design (iii) the
largest mean post-selected fidelity, $\bar F_{\mathrm{post}}=\G{noisy_control}{\NCSixFPostThree}{--}$ against
$\G{noisy_control}{\NCSixFPostOne}{--}$. As factors on the infidelity, the noise-aware gains are modest:
\G{noisy_control}{\NCSixGainRaw}{--} for (ii) and \G{noisy_control}{\NCSixGainPost}{--} for (iii) at $n=6$, and
\G{noisy_control}{\NCFourGainRaw}{--} and \G{noisy_control}{\NCFourGainPost}{--} at $n=4$. In the worst setting of
the batch the post-selected fidelity rises from \G{noisy_control}{\NCSixMinFPostOne}{--} for (i) to
\G{noisy_control}{\NCSixMinFPostThree}{--} for (iii). Figure~\ref{fig:control}(b) shows the price. The noise-aware
designs are worse than (i) at rates below the training batch (factor \G{noisy_control}{\NCSixRatioLow}{--} at
$\gamma=10^{-3}$ for (iii)) and better above it (\G{noisy_control}{\NCSixRatioHigh}{--} at $\gamma=0.1$). Post-selection
itself is far stronger than any choice of design. For the noiseless design it lowers the mean infidelity by a factor
\G{noisy_control}{\NCSixPostFactor}{--} ($n=6$) and \G{noisy_control}{\NCFourPostFactor}{--} ($n=4$), and from the best
design without post-selection, (ii), to design (iii) with it the factor is \G{noisy_control}{\NCSixGainTotal}{--} at
$n=6$. The price is the acceptance, $\bar p=\G{noisy_control}{\NCSixPThree}{--}$ for (iii) against
$\G{noisy_control}{\NCSixPTwo}{--}$ for (ii); it differs between the designs only through $T$. In short, once the
noiseless design is as fast as the controls allow, little is left for noise-aware shaping. The centre does most of the
work: the measurement of the weight removes the detected errors, and the optimiser then trades the remaining,
undetected ones against time.

\paragraph{Checks.} Every number in the table comes from a separate forward evaluation, not from the optimiser. It
integrates in physical time, passes the rates as \texttt{rates} rather than as scaled coefficients, and uses
\G{noisy_control}{\NCEvalSubsteps}{20} steps per segment. The optimiser's own values differ from it by at most
\G{noisy_control}{$\NCRecheck$}{--}, the error of the coarse time steps. The acceptances agree with the birth--death
process of Proposition~\ref{prop:noisy}(a) to \G{noisy_control}{$\NCBirthDeath$}{--}. At one point (design (iii),
mixed noise, $\gamma=0.03$, $n=\G{noisy_control}{\NCnMax}{6}$), \texttt{method='exact'} reproduces $F$,
$F_{\mathrm{post}}$ and $p$ to \G{noisy_control}{$\NCExactDiff$}{--}. A directional derivative of the objective of
(iii) agrees with central differences to a relative \G{noisy_control}{$\NCGradErr$}{--}. The whole study ($n=4$ and
$6$) took \G{noisy_control}{\NCSeconds\,s}{--} on a CPU with \G{noisy_control}{\NCThreads}{several} threads
(script \texttt{scripts/noisy\_xy\_control.py}).

\paragraph{What this shows, and what it does not.} The study shows that the library handles the whole loop of
noise-aware design: a Lindblad evolution in symmetry sectors with sparse jumps, differentiated with respect to
Hamiltonian controls and the duration; an objective defined through a central projection; and a batch of noise models,
rates and starts in a single evolution. It also gives a quantitative answer for this model. With a time-optimal
noiseless design as the reference, the noise-aware designs gain only a few per cent of the infidelity, while
post-selection gains a factor of several. Proposition~\ref{prop:noisy}(a) explains why the two noise-aware objectives
give such similar results: post-selection changes the objective only through $T$ and the weights $p_k$. The comparison
depends on the reference. With a weaker time penalty $\lambda$, the noiseless design stays well above the shortest
time, and the same optimisation then reports larger noise-aware gains, which are really gains in speed. The search is
local: there are few starts (\G{noisy_control}{\NCWarmLarge}{four} for the noise-aware designs at $n=6$), $S$ is
small, the controls are bounded and piecewise constant, and the reported designs are good solutions, not proven optima.
The sizes are small ($n=4,6$). For $n=8$ the padded $70\times70$ blocks of all nine sectors make every gradient step
much more expensive, and we did not attempt that size on a CPU. The noise models are the idealised local ones of
Section~\ref{sec:noisy}, with rates measured in units of the coupling bound. Gate-level circuits, readout errors and
non-covariant noise are not modelled, and nothing here is a claim about particular hardware.
